\documentclass[10pt,a4paper]{article}
\usepackage[margin=20mm,headheight=15pt,headsep=7mm,footskip=10mm]{geometry}
\usepackage{fontspec,amsmath,amssymb,mathtools,microtype}
\setmainfont{lmroman10-regular.otf}[BoldFont=lmroman10-bold.otf,ItalicFont=lmroman10-italic.otf,BoldItalicFont=lmroman10-bolditalic.otf]
\setsansfont{lmsans10-regular.otf}[BoldFont=lmsans10-bold.otf,ItalicFont=lmsans10-oblique.otf,BoldItalicFont=lmsans10-boldoblique.otf]
\setmonofont[Scale=.86]{lmmono10-regular.otf}[ItalicFont=lmmono10-italic.otf]
\usepackage{ragged2e}\usepackage{xcolor,booktabs,longtable,array,graphicx,enumitem,fancyhdr,needspace,titlesec,hyperref,xurl}
\definecolor{AtlasNavy}{HTML}{18334C}\definecolor{AtlasTeal}{HTML}{116D73}\definecolor{AtlasGray}{HTML}{52616D}
\hypersetup{colorlinks=true,linkcolor=AtlasTeal,urlcolor=AtlasTeal,pdfcreator={Complexity Atlas},unicode=true}
\urlstyle{same}\setlength{\parindent}{0pt}\setlength{\parskip}{4.8pt}\setlength{\emergencystretch}{3em}
\renewcommand{\baselinestretch}{1.08}\setcounter{secnumdepth}{0}\setlist{nosep,leftmargin=1.5em}
\titleformat{\section}{\Large\bfseries\sffamily\color{AtlasNavy}\raggedright}{}{0em}{}
\titleformat{\subsection}{\large\bfseries\sffamily\color{AtlasTeal}\raggedright}{}{0em}{}
\titleformat{\subsubsection}{\normalsize\bfseries\color{AtlasNavy}\raggedright}{}{0em}{}
\titlespacing*{\section}{0pt}{14pt}{6pt}\titlespacing*{\subsection}{0pt}{10pt}{5pt}
\pagestyle{fancy}\fancyhf{}\fancyhead[L]{\small\sffamily\color{AtlasGray}COMPLEXITY ATLAS}
\fancyhead[R]{\small\sffamily\color{AtlasGray}ENGLISH EDITION}\fancyfoot[C]{\small\thepage}\renewcommand{\headrulewidth}{.3pt}
\newcommand{\AtlasMath}[1]{\par\begingroup\centering\setbox0=\hbox{$\displaystyle #1$}\ifdim\wd0>1.4\linewidth\typeout{ATLAS WIDE FORMULA: \the\wd0}\fi\ifdim\wd0>\linewidth\resizebox{\linewidth}{!}{\box0}\else\box0\fi\par\endgroup}
\newcommand{\AtlasInlineMath}[1]{\begingroup\setbox0=\hbox{$#1$}\ifdim\wd0>.4\linewidth\endgroup\AtlasMath{#1}\else\endgroup$#1$\fi}
\pdfstringdefDisableCommands{\def\AtlasInlineMath#1{#1}}
\DeclareMathOperator{\polylog}{polylog}\DeclareMathOperator{\dist}{dist}\DeclareMathOperator{\diam}{diam}
\DeclareMathOperator{\rank}{rank}\DeclareMathOperator{\tr}{tr}\DeclareMathOperator{\diag}{diag}
\DeclareMathOperator{\prox}{prox}\DeclareMathOperator{\KL}{KL}\DeclareMathOperator{\Var}{Var}
\providecommand{\eps}{\varepsilon}\providecommand{\R}{\mathbb R}\providecommand{\E}{\mathbb E}
\providecommand{\mathbbm}[1]{\mathbf{#1}}\providecommand{\mathds}[1]{\mathbf{#1}}
\providecommand{\llbracket}{[\![}\providecommand{\rrbracket}{]\!]}
\begin{document}
\raggedright

{\LARGE\bfseries\sffamily\color{AtlasNavy} Gradient-norm minimization for convex functions: non-Euclidean Hölder smoothness; exact function values and gradients\par}\vspace{5pt}
\textbf{Setting ID.} \texttt{SO\_\allowbreak{}FO\_\allowbreak{}C07\_\allowbreak{}holder\_\allowbreak{}FG}

Single-objective optimization

Find a point satisfying the stationarity criterion specified under the accuracy measure.

\textbf{Status.} Bounds missing

\textbf{TASK tags.} formulation: First order stationarity; task family: Single objective

\textbf{MODEL tags.} convexity: Convex; geometry: Full space; geometry: Norm ball; smoothness: Holder gradient; smoothness: Holder smooth

\textbf{ORACLE tags.} access order: First-order interface; feedback quality: Exact; interface family: Fg; operation: Value gradient; public information: Computable access; public information: Hidden objects explicitly declared; public information: Known upper bound; public information: Known exact value; public information: Public inputs explicitly declared

\textbf{ALGORITHMS tags.} adaptivity: Adaptive; computation: Internal arithmetic uncharged; randomization: Private randomness allowed; restriction: No additional span restriction explicitly imposed

\textbf{OUTPUT tags.} certificate: No additional output certificate required; feasibility: Required; object: Point; postprocessing: Unqueried output explicitly allowed

\textbf{METRIC tags.} criterion: First order stationarity; criterion: Gradient norm

\textbf{GUARANTEE tags.} aggregation: Per instance; mode: Expectation; randomness: Algorithm private seed

\textbf{RESOURCE tags.} budget: Fixed worst-case budget; measure: Oracle calls; precision: Bit precision not counted

\needspace{7\baselineskip}\section{1. Setting specification}
\needspace{4\baselineskip}\subsection{Task}
Find a point satisfying the stationarity criterion specified under the accuracy measure.


\AtlasMath{\min_{x\in \mathbb{R}^d}f(x)}


\needspace{4\baselineskip}\subsection{Admissible model}
\AtlasInlineMath{f} is globally convex and satisfies \AtlasInlineMath{\|\nabla f(x)-\nabla f(y)\|_{p^\star}\le\beta_\nu\|x-y\|_p^\nu}. A global minimizer satisfies \AtlasInlineMath{\|x^\star-x_0\|_p\le R}, where \AtlasInlineMath{1<p<\infty} and \AtlasInlineMath{0<\nu\le1}.

Admissible parameter domain:


\AtlasMath{d\in\mathbb{N}_{\ge1},\quad \varepsilon>0,\quad R>0,\quad \beta_\nu>0,\quad \nu>0,\quad \nu\le1,\quad p\ge1,\quad p>1}


\needspace{4\baselineskip}\subsection{Oracle interface}
Public numerical inputs: \AtlasInlineMath{d,\ p,\ \nu,\ \varepsilon,\ x_0}.

Public upper bounds: \AtlasInlineMath{\beta_\nu,\ R}.

The unknown object is the objective \AtlasInlineMath{f}.

Query points lie in \AtlasInlineMath{\mathbb{R}^d}; feedback consists of exact function values and gradients:


\AtlasMath{(f(x),\nabla f(x))}


All responses are consistent with one fixed function instance. Exact values and derivatives agree across repeated queries at the same point.

\needspace{4\baselineskip}\subsection{Output}
Return a point \AtlasInlineMath{\widehat x\in\mathbb R^d}.

The output may be unqueried or constructed from the query history, provided it is feasible. No additional function value or accuracy certificate is required.

\needspace{4\baselineskip}\subsection{Loss or accuracy criterion}

\AtlasMath{\|\nabla f(\widehat{x})\|_{p^\star}}


\AtlasInlineMath{\varepsilon} is a gradient-norm tolerance. Expected norm and expected squared norm are different performance requirements.

\needspace{4\baselineskip}\subsection{Resource accounting}
One declared full first-order response costs one call; its vector coordinates are not charged separately.

Cost is bounded uniformly on every execution path. All calls used for initialization, repetition, batching, or candidate verification are counted.

Real arithmetic, precision, and storage applied to already available information are uncharged. Public geometric operations are limited to those explicitly granted by the oracle. Query complexity alone does not imply arithmetic, memory, or communication complexity.

\needspace{4\baselineskip}\subsection{Algorithm views and success guarantees}
Private randomization is allowed. Queries, stopping decisions, and postprocessing may adapt to observed feedback.

No additional span, zero-respecting, linear-iteration, or memory restriction is imposed.

Interaction is adaptive to the feedback history. Every actual response is charged separately, including responses in parallel or batch queries. Internal real arithmetic, storage, and postprocessing are uncharged.

Select an algorithm using the public parameters, and then require the following for every admissible instance and response rule:


\AtlasMath{\exists A\ \forall I,\mathcal{O}:\ \mathbb{E}_A[\ell(I,\widehat{x})]\le\varepsilon}


The probability space includes the algorithm's private randomness. The instance and its admissible oracle law are fixed before randomness is averaged.

\needspace{7\baselineskip}\section{2. Notation and parameters}
\AtlasInlineMath{x,\widehat{x}}: Query points and final output; their permitted spaces are specified separately by the oracle and output contracts \textbf{Domain:} \AtlasInlineMath{\widehat{x}\in \mathbb{R}^d}.

\AtlasInlineMath{f}: Unknown objective, with domain and structure specified by the model

\AtlasInlineMath{\nabla f}: Full gradient of the differentiable objective or smooth term with respect to the decision variable

\AtlasInlineMath{x_0}: Public initial point satisfying the stated initial-distance or initial-suboptimality condition

\AtlasInlineMath{f^\star}: Optimal value of the complete task; the oracle need not return it


\AtlasMath{f^\star=\inf_{x\in \mathbb{R}^d}f(x)}


\AtlasInlineMath{x^\star}: A global minimizer whose existence is promised by the model; unique under strong convexity

\AtlasInlineMath{\mathbb{R}^d}: Real vector space with \AtlasInlineMath{d} coordinates

\AtlasInlineMath{\|\cdot\|_2,\langle\cdot,\cdot\rangle}: Euclidean norm and inner product

\AtlasInlineMath{\|\cdot\|_p}: Standard finite-dimensional norm with the indicated exponent \textbf{Domain:} Endpoint norm \AtlasInlineMath{\|x\|_\infty=\max_{1\le j\le d}|x_j|}.


\AtlasMath{\|x\|_a=(\sum_{j=1}^d|x_j|^a)^{1/a},\quad1\le a<\infty}


\AtlasInlineMath{d}: Dimension; a public positive integer whose dependence is tracked jointly \textbf{Domain:} \AtlasInlineMath{d\in\mathbb{N}_{\ge1}}.

\AtlasInlineMath{\beta}: Public upper bound on the gradient Lipschitz or smoothness constant \textbf{Domain:} \AtlasInlineMath{\beta>0}.

\AtlasInlineMath{\beta_\nu}: Public upper bound on the gradient Hölder constant \textbf{Domain:} \AtlasInlineMath{\beta_\nu>0}.

\AtlasInlineMath{R}: Upper bound on the distance from the public initial point to a minimizer \textbf{Domain:} \AtlasInlineMath{R>0}.

\AtlasInlineMath{\varepsilon}: Target tolerance for the underlying error defined by the accuracy measure \textbf{Domain:} \AtlasInlineMath{\varepsilon>0}.

\AtlasInlineMath{p}: Public norm exponent specifying the feasible-set or localization geometry \textbf{Domain:} \AtlasInlineMath{1\le p\le\infty}.

\AtlasInlineMath{\nu}: Hölder exponent; public, with fixed or varying status stated in each comparison \textbf{Domain:} \AtlasInlineMath{0<\nu\le1}.

\AtlasInlineMath{C}: Sufficiently large universal positive constant, not specified numerically and potentially different between theorems. \textbf{Domain:} \AtlasInlineMath{C>1}.

\AtlasInlineMath{r_p}: Higher-order uniform-convexity exponent \AtlasInlineMath{r_p=\max\{2,p\}}.


\AtlasMath{r_p=\max\{2,p\}}


\AtlasInlineMath{c_{p,\nu}}: Positive multiplicative constant, not specified numerically, depending only on the following parameters fixed in advance: \AtlasInlineMath{p},\AtlasInlineMath{\nu}. \textbf{Domain:} \AtlasInlineMath{c_{p,\nu}>0}.

\AtlasInlineMath{C_{p,\nu}}: Positive multiplicative constant, not specified numerically, depending only on the following parameters fixed in advance: \AtlasInlineMath{p},\AtlasInlineMath{\nu}. \textbf{Domain:} \AtlasInlineMath{C_{p,\nu}>0}.

\AtlasInlineMath{C_{\nu}}: Positive multiplicative constant, not specified numerically, depending only on the following parameters fixed in advance: \AtlasInlineMath{\nu}. \textbf{Domain:} \AtlasInlineMath{C_{\nu}>0}.

\AtlasInlineMath{p^\star}: Conjugate exponent of \AtlasInlineMath{p}


\AtlasMath{1/p+1/p^\star=1}


\AtlasInlineMath{\ell}: The underlying loss or error defined by the accuracy measure; its expected, probabilistic, or worst-case aggregation is specified by the success contract

\AtlasInlineMath{I,\mathcal{O},A}: Admissible instance, a response rule consistent with that instance, and an algorithm; the algorithm is selected before the unknown instance

\AtlasInlineMath{\mathbb{E}}: Expectation over the declared randomness after fixing the instance and admissible response law

\AtlasInlineMath{O,\Omega,\Theta,\widetilde{O},\widetilde\Omega}: Notation for asymptotic upper and lower bounds, constant-factor matching rates, and explicitly permitted logarithmic factors. Hidden-constant dependence is specified in each regime.

\AtlasInlineMath{\log}: Natural logarithm; changing to another fixed base affects only a constant

\needspace{12\baselineskip}\section{3. Complexity bracket}
\needspace{4\baselineskip}\subsection{Compared quantities}
\needspace{4\baselineskip}\subsubsection{Q1 - target accuracy cost}
\AtlasMath{C^\star_{v,\pi}(\varepsilon;\theta)=\inf\{B\ge0:\exists\mathcal A\in\mathfrak A_v,\ C_\pi(\mathcal A,I)\le B\ \text{and}\ G_v(\mathcal A,I;\varepsilon)\ \forall I\}}
\textbf{Definition.} One declared full first-order response costs one call; its vector coordinates are not charged separately.

\textbf{Scope.} declared\_joint\_parameter\_window

\textbf{Optimization direction.} minimize

Here v is the selected algorithm view, G\_v its complete success guarantee, I an admissible fixed instance, theta the remaining parameters, and C\_pi the declared resource projection. C\_pi uses the view's fixed-path or expected-work budget convention, including initialization. No change of guarantee or scalarization is implicit.

\needspace{4\baselineskip}\subsection{R1 - Nontrivial Hölder regime in a general norm}
\textbf{Region.} \AtlasInlineMath{1<p<\infty,\ 0<\nu\le1,\quad\varepsilon\le c_{p,\nu}\beta_\nu R^\nu}

\textbf{Quantity.} Q1

\textbf{Resource.} One declared full first-order response costs one call; its vector coordinates are not charged separately.

\textbf{Status.} Lower bound missing

\textbf{Lower bound.} No admitted bound for this comparison

\textbf{Upper bound.} B1, B2

\AtlasMath{\min\{C_{p,\nu}(\beta_\nu R^\nu/\varepsilon)^{r_p/[1+(r_p+1)\nu]},C_\nu d\log(2+\beta_\nu R^\nu/\varepsilon)\}}
\textbf{Comparison argument.} The comparison uses the same model, query domain, output requirement, success guarantee, and resource convention. The upper envelope minimizes over applicable algorithms; the lower envelope maximizes over applicable hard instances. A finite sum of lower bounds is used only up to a constant factor. Rate comparison: The upper bound is verified for general p. The sufficient high-dimensional condition of the randomized lower bound is not fully numerical; this does not mean that no lower bound exists. Ratio: None

\needspace{7\baselineskip}\section{4. Bounds and supporting evidence}
The following conditions apply to every bound in this profile, in addition to each bound's stated region and algorithm scope.

\textbf{Conditions.} 
\AtlasMath{1<p<\infty,\ 0<\nu\le1,\quad\varepsilon\le c_{p,\nu}\beta_\nu R^\nu}


\textbf{Dimension conditions.} \AtlasInlineMath{1<p<\infty,\ 0<\nu\le1,\quad\varepsilon\le c_{p,\nu}\beta_\nu R^\nu}

\textbf{Oracle mapping.} Applicable interfaces: exact function values and gradients, exact gradients. Algorithms may use only these interfaces and the auxiliary operations explicitly declared public.

\textbf{Guarantee mapping.} A deterministic error bound for every instance also gives the same expected-error bound. Probability guarantees are used only when an explicit conversion to the target contract is recorded.

\textbf{Output mapping.} The output must lie in the declared output space and satisfy its feasibility, support, or atomic-hull restrictions. Permission to return an unqueried point does not remove these restrictions.

\textbf{Resource mapping.} Every primitive response is charged. Sample, component, proximal, and communication costs use the units declared for the target setting. Resource coordinates achieved by one algorithm remain a joint vector; coordinatewise minima from different algorithms are not combined.

\needspace{4\baselineskip}\subsection{B1 - Gradient-norm minimization under non-Euclidean Hölder smoothness}
\textbf{Admission.} Included in the bracket

\textbf{Direction.} upper

\textbf{Witness or lower-bound scope.} Deterministic witness

\AtlasMath{C_{p,\nu}(\beta_\nu R^\nu/\varepsilon)^{r_p/[1+(r_p+1)\nu]},\quad r_p=\max\{2,p\}}
\textbf{Sources.} \href{https://arxiv.org/pdf/2609.01122}{[1]}

\textbf{Source location.} \href{https://arxiv.org/pdf/2609.01122}{[1]} theorem: Theorem 5.3, PDF p. 6; Assumption 4.1; the public auxiliary minimization in equation (1)

\needspace{4\baselineskip}\subsection{B2 - Localized cutting planes for Hölder-smooth functions}
\textbf{Admission.} Included in the bracket

\textbf{Direction.} upper

\textbf{Witness or lower-bound scope.} Deterministic witness

\AtlasMath{C_\nu d\log(2+\beta_\nu R^\nu/\varepsilon)}
\textbf{Sources.} \href{https://arxiv.org/pdf/2101.09740v2}{[2]} \href{https://arxiv.org/pdf/2609.33360v2}{[3]} \href{https://arxiv.org/pdf/2609.10874}{[4]} \href{https://arxiv.org/pdf/1405.4980}{[5]} \href{https://www2.isye.gatech.edu/~nemirovs/MOR_AccuracyCertificates.pdf}{[6]} \href{https://arxiv.org/pdf/2609.20701}{[7]}

\textbf{Local argument.} A1, A2, A3

\needspace{7\baselineskip}\section{5. Local arguments}
\needspace{4\baselineskip}\subsection{A1 - Smooth localization and gradient-norm conversion}
\textbf{Scope.} Gradient access and a public ball localizing a global minimizer; public regularization when only initial gap is given.

\textbf{Algorithm scope.} Deterministic upper witnesses; imported lower theorems retain their own deterministic or randomized scopes.

If a global minimizer lies within distance \AtlasInlineMath{R} of \AtlasInlineMath{x_0}, query on \AtlasInlineMath{B(x_0,R)}. Since \AtlasInlineMath{\nabla f(x^\star)=0}, gradients there have norm at most \AtlasInlineMath{2\beta R}. The centroid certificate gives \AtlasInlineMath{O(d\log(2+\beta R^2/\epsilon))} gradient calls. Smoothness gives \AtlasInlineMath{f(x_0)-f^\star\le\beta R^2/2}, including the zero-query endpoint.

With global \AtlasInlineMath{\mu}-strong convexity and initial gap \AtlasInlineMath{\Delta}, use radius \AtlasInlineMath{\sqrt{2\Delta/\mu}}, obtaining \AtlasInlineMath{O(d\log(2+\kappa\Delta/\epsilon))}. At \AtlasInlineMath{\mu=\beta}, equality in strong monotonicity and Lipschitz continuity forces \AtlasInlineMath{\nabla f(x)-\nabla f(y)=\beta(x-y)}. One gradient identifies the minimizer of this known-curvature quadratic plus unknown affine term. Opposite translations exclude zero calls at nontrivial accuracy.

Under global Hölder smoothness in \AtlasInlineMath{\ell_q}, convert to Euclidean quantities
\AtlasInlineMath{\beta_E=\beta_\nu d^{(1+\nu)(1/q-1/2)_+}} and
\AtlasInlineMath{R_E=Rd^{(1/2-1/p)_+}}. Gradients on the localization ball are at most \AtlasInlineMath{\beta_E(2R_E)^\nu}, giving
\AtlasInlineMath{C_\nu d\log(2+\beta_E R_E^{1+\nu}/\epsilon)}.
A boundary minimizer of an arbitrary constrained problem does not justify the zero-gradient premise. With two public norm-distance bounds, their intersection lies in the Euclidean ball of radius \AtlasInlineMath{R_{2,+}}, while the Euclidean ball of radius \AtlasInlineMath{R_{2,-}} satisfies both promises. The former gives a certificate upper bound; embedding the one-dimensional hard family in the latter gives the logarithmic lower bound.

For gradient-norm accuracy, use \AtlasInlineMath{\|\nabla f(x)\|^2\le2\beta(f(x)-f^\star)}. Under initial distance, target gap \AtlasInlineMath{\epsilon^2/(2\beta)}. Under initial gap alone, instead set

\AtlasMath{F_\lambda=f+\lambda\|x-x_0\|^2/2,\qquad\lambda=\epsilon^2/(32\Delta).}

Its minimizer lies within \AtlasInlineMath{8\Delta/\epsilon}. A certificate at gap \AtlasInlineMath{\epsilon^2/(128\beta)} gives \AtlasInlineMath{\|\nabla F_\lambda(\widehat x)\|\le\epsilon/4}, while the public regularizer contributes at most \AtlasInlineMath{\epsilon/4}. Each regularized gradient costs one original gradient, and the budget is \AtlasInlineMath{O(d\log(2+\beta\Delta/\epsilon^2))}.

For lower transfers, strong convexity gives \AtlasInlineMath{f-f^\star\le\|\nabla f\|^2/(2\mu)}. Choose the hard subclass modulus proportional to \AtlasInlineMath{\epsilon/R}, or \AtlasInlineMath{\epsilon^2/\Delta}, for the two respective initializations. Gradient accuracy forces constant relative gap reduction. Apply the cited deterministic chain or randomized matrix-query theorem in its original dimension window; truncate only to constructions that fit. This preserves initialization, guarantee, and algorithm scope.

For an already strongly convex model, the gap threshold required for the lower transfer is \AtlasInlineMath{\epsilon^2/(2\mu)}, whereas the sufficient gap for the upper transfer is \AtlasInlineMath{\epsilon^2/(2\beta)}. Retain these different logarithm arguments rather than silently canceling their condition-number shift. For \AtlasInlineMath{1<\kappa\le4}, gradient descent and the nested one-dimensional construction yield the stated comparable near-unit denominators. At \AtlasInlineMath{\kappa=1}, the known-curvature quadratic is solved in one call.

For matching \AtlasInlineMath{\ell_p} Hölder geometry, the descent model along a dual steepest-descent direction gives
\AtlasInlineMath{f(x)-f^\star\ge c_\nu\beta_\nu^{-1/\nu}\|\nabla f(x)\|_{p^\star}^{(1+\nu)/\nu}}.
Targeting gap \AtlasInlineMath{c_\nu\beta_\nu^{-1/\nu}\epsilon^{(1+\nu)/\nu}} in the localization certificate therefore gives
\AtlasInlineMath{C_\nu d\log(2+\beta_\nu R^\nu/\epsilon)}.

\textbf{Limits of the argument.} An arbitrary constrained minimizer may be on the boundary and need not have zero gradient. Fixed-dimensional lower transfers retain theorem dimension windows.

\textbf{Sources.} \href{https://arxiv.org/pdf/2101.09740v2}{[2]} \href{https://arxiv.org/pdf/2609.33360v2}{[3]} \href{https://arxiv.org/pdf/2609.10874}{[4]}

\needspace{4\baselineskip}\subsection{A2 - Descent, PL contraction, and zero-query endpoints}
\textbf{Scope.} The exact smooth, PL, Clarke, Moreau, or composite residual specified in the associated bound.

\textbf{Algorithm scope.} Deterministic upper witnesses and elementary indistinguishability lower bounds.

Global \AtlasInlineMath{\beta}-smoothness and boundedness below imply

\AtlasMath{f(x-\nabla f(x)/\beta)\le f(x)-\|\nabla f(x)\|^2/(2\beta).}

Thus \AtlasInlineMath{\|\nabla f(x_0)\|^2\le2\beta\Delta}. Summing the decrease and returning the queried point with least gradient norm gives \AtlasInlineMath{O(1+\beta\Delta/\epsilon^2)} calls; values are unnecessary. If the initial gradient bound already meets the tolerance, output \AtlasInlineMath{x_0} without querying.

A globally \AtlasInlineMath{\beta}-Lipschitz gradient gives Hessian norm at most \AtlasInlineMath{\beta}, so the second-order zero-query regime additionally requires permitted negative curvature at least \AtlasInlineMath{\beta}.

For a global \AtlasInlineMath{\mu}-PL objective,
\AtlasInlineMath{\mathrm{gap}_{t+1}\le(1-1/\kappa)\mathrm{gap}_t}.
This proves the exact ceiling formula for \AtlasInlineMath{1<\kappa}. At \AtlasInlineMath{\kappa=1}, one step is exact, and opposite translated quadratics exclude zero calls at nontrivial accuracy. At \AtlasInlineMath{\epsilon\ge\Delta}, the initial point suffices.

In one dimension, PL excludes nonoptimal stationary points. The minimizer set is an interval. Integrating
\AtlasInlineMath{|f'|\ge\sqrt{2\mu(f-f^\star)}}
along the decreasing side shows that a minimizer lies within \AtlasInlineMath{\sqrt{2\Delta/\mu}} of \AtlasInlineMath{x_0}. Derivative-sign bisection and smoothness therefore give \AtlasInlineMath{O(\log(2+\kappa\Delta/\epsilon))} calls.

Clarke subgradients of a globally \AtlasInlineMath{L}-Lipschitz function have norm at most \AtlasInlineMath{L}, so Goldstein residuals are at most \AtlasInlineMath{L}, also at radius zero. For a \AtlasInlineMath{\rho}-weakly convex function, proximal optimality gives
\AtlasInlineMath{\|\nabla f_{1/(2\rho)}(x_0)\|\le L}; the proximal objective inequality gives its square at most \AtlasInlineMath{4\rho\Delta}.
For a smooth composite objective, proximal-gradient descent gives
\AtlasInlineMath{F(y)\le F(x)-\|\mathcal G_\beta(x)\|^2/(2\beta)},
so \AtlasInlineMath{\|\mathcal G_\beta(x_0)\|^2\le2\beta\Delta}.
These bounds certify the corresponding zero-query outputs; the output contract does not require computing the residual.

\textbf{Limits of the argument.} Each endpoint uses its own residual and model; the metrics are not interchangeable.

\needspace{4\baselineskip}\subsection{A3 - Subgradient certificates and centroid cuts}
\textbf{Scope.} One full subgradient is sufficient. Public compact convex geometry and exact internal geometric computation are available.

\textbf{Algorithm scope.} Deterministic upper-bound witness; valid as an upper bound for randomized views.

Let \AtlasInlineMath{Q\subset\mathbb R^d} be public, compact, convex, and full dimensional, with diameter at most \AtlasInlineMath{D} in the stated norm. Suppose all returned subgradients have dual norm at most \AtlasInlineMath{L}. Query \AtlasInlineMath{g_i\in\partial f(x_i)} at the centroid of the current body and retain \AtlasInlineMath{\langle g_i,x_i-x\rangle\ge0}. The centroid halfspace theorem bounds the retained volume by \AtlasInlineMath{(1-e^{-1})^T\operatorname{vol}(Q)}.

The computable certificate is

\AtlasMath{\Gamma_T=\min_{\lambda\in\Delta_T}\max_{x\in Q}\sum_i\lambda_i\langle g_i,x_i-x\rangle
=\max_{x\in Q}\min_i\langle g_i,x_i-x\rangle.}

The equality is convex minimax. Subgradient inequalities and convexity imply
\AtlasInlineMath{f(\sum_i\lambda_ix_i)-\min_Qf\le\Gamma_T}.
If \AtlasInlineMath{\Gamma_T>\epsilon}, choose \AtlasInlineMath{z\in Q} with all cut margins greater than \AtlasInlineMath{\epsilon}. The copy \AtlasInlineMath{(1-a)z+aQ}, where \AtlasInlineMath{a=\epsilon/[2(\epsilon+LD)]}, survives every cut. Its volume is \AtlasInlineMath{a^d\operatorname{vol}(Q)}, contradicting centroid shrinkage once \AtlasInlineMath{T>C d\log(2+LD/\epsilon)}. If the body loses volume earlier, the same positive-volume argument already certifies the error. No function values or inscribed-ball estimate are required.

For \AtlasInlineMath{Q=RB_p^d} and \AtlasInlineMath{\ell_q}-Lipschitz objectives, use \AtlasInlineMath{D=2Rd^{(1/q-1/p)_+}}. Returning the center has error at most \AtlasInlineMath{LRd^{(1/q-1/p)_+}}, giving the coarse zero-query regime. Exact centroid, minimax, and public geometric calculations are uncharged internal operations.

For a deterministic selector with total movement at most \AtlasInlineMath{W(T)} on nested convex sets, instead impose deep cuts of margin \AtlasInlineMath{2LW(T)/T}. If all \AtlasInlineMath{T} sets are nonempty, every move has length at least \AtlasInlineMath{2W(T)/T}, contradicting the movement bound. At the first empty intersection, minimax gives error at most \AtlasInlineMath{2LW(T)/T}. Substitution of the cited stable-movement theorem yields the recorded nondual upper bound. Its deterministic information scope is distinct from the scope of randomized implementations.

\textbf{Limits of the argument.} The stable-selector branch imports its movement theorem and does not claim efficient deterministic arithmetic time.

\textbf{Sources.} \href{https://arxiv.org/pdf/1405.4980}{[5]} \href{https://www2.isye.gatech.edu/~nemirovs/MOR_AccuracyCertificates.pdf}{[6]} \href{https://arxiv.org/pdf/2609.20701}{[7]}

\needspace{7\baselineskip}\section{References}
[1] Nico Pelleriti; Maryam Shiran; David Martínez-Rubio; Max Zimmer; Sebastian Pokutta. Optimal Gradient-Norm Minimization in Non-Euclidean Hölder-Smooth Convex Optimization. 2026-09-01; preprintv1.

\url{https://arxiv.org/pdf/2609.01122}

[2] Yoel Drori; Adrien Taylor. On the oracle complexity of smooth strongly convex minimization. 2021; arXiv:2101.09740v2.

\url{https://arxiv.org/pdf/2101.09740v2}

[3] Yoel Drori. Exact oracle complexity for function-value suboptimality under weighted initial conditions. 2026-10-16; preprintv2.

\url{https://arxiv.org/pdf/2609.33360v2}

[4] Ohad Shamir. A Simple Complexity Lower Bound for Solving Ax=b. 2026-09-09; preprintv1.

\url{https://arxiv.org/pdf/2609.10874}

[5] Sébastien Bubeck. Convex Optimization: Algorithms and Complexity. 2015.

\url{https://arxiv.org/pdf/1405.4980}

[6] Arkadi Nemirovski; Shmuel Onn; Uriel G. Rothblum. Accuracy Certificates for Computational Problems with Convex Structure. 2010.

\url{https://www2.isye.gatech.edu/~nemirovs/MOR_AccuracyCertificates.pdf}

[7] David Martínez-Rubio; Cristóbal Guzmán. Stable Movement for Nondual Lipschitz Convex Optimization: Efficiency and Nearly Optimal Oracle Rates. 2026-09-17; preprintv1.

\url{https://arxiv.org/pdf/2609.20701}

\end{document}
